\question

Realice cada uno de los ejemplos de\href{https://drive.google.com/file/d/1OFmqAdh1-XTbVeuHS5Mn-c7aSAErCQe2/view}{\textcolor{red}{\faIcon{file-pdf}}\mintinline{latex}|t1.pdf|}:\href{https://drive.google.com/file/d/1rbNvSLlXMA1YnocPLROx05F6oJ-7Fvr6/view}{\textcolor{violet}{\faIcon{file}}\mintinline{latex}|helloworld.tex|},
\href{https://drive.google.com/file/d/1nL1XyMtPDasrUvjSZe6VIK8f1OXC1M9K/view}{\textcolor{violet}{\faIcon{file}}\mintinline{latex}|article.tex|},\href{https://drive.google.com/file/d/1s7p8jGuN_j0TmOr0vXCmb4ENJQFdZuhz/view}{\textcolor{violet}{\faIcon{file}}\mintinline{latex}|book.tex|},\\\href{https://drive.google.com/file/d/1FwOW64Jd24TQpoHid0SUy4QbpD6IntSU/view}{\textcolor{violet}{\faIcon{file}}\mintinline{latex}|report.tex|},\href{https://drive.google.com/file/d/1CcgX8S3e-8eVyT29JlFR4L6X_LVgXl9w/view}{\textcolor{violet}{\faIcon{file}}\mintinline{latex}|exam.tex|},\href{https://drive.google.com/file/d/1G0sE_5IZ3Di322UcV4XJAzcyRFH4tbyC/view}{\textcolor{violet}{\faIcon{file}}\mintinline{latex}|beamer.tex|} y\href{https://drive.google.com/file/d/1ACXGEDKyuMb4JJr4pZ_dlAeLPUnh9zbB/view}{\textcolor{violet}{\faIcon{file}}\mintinline{latex}|poster.tex|}.
En cada archivo, modifique\\\mintinline{latex}|\title{<Título>}| y \mintinline{latex}|\author{<Nombre>}|,
siempre que existan.

\question

Digite el ejercicio final de\href{https://drive.google.com/file/d/1OFmqAdh1-XTbVeuHS5Mn-c7aSAErCQe2/view}{\textcolor{red}{\faIcon{file-pdf}}\mintinline{latex}|t1.pdf|},
en el preámbulo incluya \mintinline{latex}|\documentclass{article}|.

\question

Importe \mintinline{latex}|\usepackage[shortlabels]{enumitem}|\footnote{Estudie la segunda sección de~\url{http://mirrors.ctan.org/macros/latex/contrib/enumitem/enumitem.pdf}}
y realice el ejemplo de listas numeradas de\href{https://drive.google.com/file/d/1O9foAwQ7XXrF28VHOq9FNcEt7kRHy3t5/view}{\textcolor{red}{\faIcon{file-pdf}}\mintinline{latex}|t2.pdf|}
con el entorno \mintinline{latex}|\begin{enumerate}[(i)]|.
¿Qué observa si crea las listas numeradas con el entorno
\mintinline{latex}|\begin{enumerate}[label=(\arabic*)]|?

\question

Importe \mintinline{latex}|\usepackage{siunitx}|\footnote{Estudie la tercera sección de~\url{http://mirrors.ctan.org/macros/latex/contrib/siunitx/siunitx.pdf}}
y cree una tabla de las magnitudes físicas fundamentales y otra tabla
de las magnitudes físicas derivadas, según el Sistema
Internacional\footnote{Vea~\url{https://en.wikipedia.org/wiki/International_System_of_Quantities}},
que indique en cada columna su nombre de la magnitud, su símbolo y su descripción.

\question

Importe \mintinline{latex}|\usepackage{xcolor}| y realice el espectro de
colores\footnote{Vea la tercera sección de \url{http://mirrors.ctan.org/macros/latex/contrib/xcolor/xcolor.pdf}}.

\begin{figure}[htbp]
	\centering
	\small
	\newcount\WL \unitlength.75pt
	\begin{picture}(460,60)(355,-10)
		\sffamily \tiny \linethickness{1.25\unitlength} \WL=360
		\multiput(360,0)(1,0){456}%
		{{\color[wave]{\the\WL}\line(0,1){50}}\global\advance\WL1}
		\linethickness{0.25\unitlength}\WL=360
		\multiput(360,0)(20,0){23}%
		{\picture(0,0)
			\line(0,-1){5} \multiput(5,0)(5,0){3}{\line(0,-1){2.5}}
			\put(0,-10){\makebox(0,0){\the\WL}}\global\advance\WL20
			\endpicture}
	\end{picture}
\end{figure}

\question

Digite el ejercicio final de\href{https://drive.google.com/file/d/1O9foAwQ7XXrF28VHOq9FNcEt7kRHy3t5/view}{\textcolor{red}{\faIcon{file-pdf}}\mintinline{latex}|t2.pdf|},
en el preámbulo incluya \mintinline{latex}|\documentclass{standalone}|.

\question

Digite la sección {\sffamily\color{AmurmapleBlue!80!black}1.6 The Gauss Transformation}
del\href{https://www.ceremade.dauphine.fr/~fejoz/Enseignement/ds2021/Brin-Stuck_2002_Introduction-to-Dynamical-Systems.pdf}{\textcolor{red}{\faIcon{file-pdf}}texto},
sin incluir los ejercicios.

\question

Importe \mintinline{latex}|\usepackage{diffcoeff}| y digite la
sección {\sffamily\color{AmurmapleBlue!80!black}Transformación de Hopf-Cole}
del\href{https://weblibrary.mila.edu.my/upload/ebook/engineering/2014_Book_IntroductionToPartialDifferent.pdf#page=338}{\textcolor{red}{\faIcon{file-pdf}}texto}
para transformar la ecuación potencial de Burgers en la ecuación de
difusión lineal.
\begin{equation*}
	\boxed{
		\diffp{v}{t}=
		\gamma\diffp[2]{v}{x}\quad
		\xRightarrow{v\left(t,x\right)=\exp\left(\alpha\varphi\left(t,x\right)\right)}\quad
		\diffp{\varphi}{t}=
		\gamma\diffp[2]{\varphi}{x}+
		\gamma\alpha{\left(\diffp{\varphi}{x}\right)}^{2}.
	}
\end{equation*}

\question

Importe \mintinline{latex}|\usepackage{subfig}| y realice el
siguiente mosaico de figuras (referenciales).

\begin{figure*}
	\captionsetup[figure]{margin=10pt}
	\subfloat[Primero izquierda.]{\includegraphics[width=.2\paperwidth]{example-image-duck}}
	\hfill
	\subfloat[Primero derecha.]{\includegraphics[width=.2\paperwidth]{example-image-duck}}\\[-10pt]
	\caption{Primera figura.}

	\subfloat[Segundo izquierda.]{\includegraphics[width=.2\paperwidth]{example-image-duck}}
	\hfill
	\subfloat[Segundo derecha.]{\includegraphics[width=.2\paperwidth]{example-image-duck}}\\[-10pt]
	\caption{Segunda figura.}

	\subfloat[Tercero izquierda.]{\includegraphics[width=.2\paperwidth]{example-image-duck}}
	\hfill
	\subfloat[Tercero derecha.]{\includegraphics[width=.2\paperwidth]{example-image-duck}}\\[-10pt]
	\caption{Tercera figura.}
\end{figure*}

\question

Importe \mintinline{latex}|\usepackage{optidef}| y digite la sección
{\sffamily\color{AmurmapleBlue!80!black}1.8}
del\href{http://cnsa.cmlaboratory.com/pdf/monograf/Andreasson,%20Evgrafov,%20Patriksson.%20An%20introduction%20to%20continuous%20optimization.pdf#page=28}{\textcolor{red}{\faIcon{file-pdf}}texto}.

\question

Importe \mintinline{latex}|\usepackage{amsthm}| y digite el teorema
	{\sffamily\color{AmurmapleBlue!80!black}1.4}
del\href{https://www.davcollegekanpur.ac.in/assets/ebooks/Maths/Hungerford-%20Algebra%20Graduate%20Texts%20in%20Mathematics.pdf}{\textcolor{red}{\faIcon{file-pdf}}texto}.

\question

Cree una presentación de al menos cinco \mintinline{latex}|frame|s
acerca del algoritmo $QR$ con la plantilla
\url{https://github.com/Numerical-Analysis-2024/beamertheme-numeric}.
Agregue sus propias referencias en \mintinline{latex}|references.bib|.
